% ID: SEL-D0056-OB-E01
% ESTADO: validada. S057.
Sean $x,y$ las unidades semanales de A y B. Maximizamos $70x+50y$ con
\[
2x+4y\le100,\qquad5x+3y\le110,\qquad x,y\ge0.
\]
Las dos fronteras de capacidad se cortan en $(10,20)$. En los ejes, las restricciones más exigentes dan $(22,0)$ y $(0,25)$. Junto con el origen forman el recinto.
\begin{center}
\setlength{\unitlength}{0.9mm}
\begin{picture}(92,62)(-10,-7)
\qbezier(0,0)(38,0)(76,0)
\qbezier(0,0)(0,26)(0,52)
\put(77,0){$x$}\put(2,52){$y$}
\qbezier(0,50)(15,45)(30,40)
\qbezier(30,40)(48,20)(66,0)
\put(-9,51){\scriptsize $(0,25)$}
\put(31,41){\scriptsize $(10,20)$}
\put(61,-5){\scriptsize $(22,0)$}
\end{picture}
\end{center}
\[
\begin{array}{c|rrrr}
(x,y)&(0,0)&(22,0)&(10,20)&(0,25)\\\hline
\text{Beneficio}&0&1540&1700&1250
\end{array}
\]
El mayor beneficio se alcanza en el vértice $(10,20)$. Como las coordenadas son enteras, la solución también es óptima si las unidades son indivisibles. Se deben producir $10$ unidades de A y $20$ de B, con beneficio de $\boxed{1700\text{ euros}}$.
